% CAPÍTULO 1 — INTRODUÇÃO -----------------------------------------------------------------------

\chapter{Introdução}
\label{cap:introducao}

Poucos setores da economia brasileira se transformaram tão rápido quanto o de Tecnologia da Informação (TI). A transformação digital de empresas e órgãos públicos tem puxado a demanda por profissionais especializados \cite{Brasscom2024, Fecomercio2024}, e entender como esse mercado se comporta passou a ser uma preocupação tanto de quem formula política pública quanto de gestores e dos próprios trabalhadores.

No Brasil, as duas principais fontes oficiais sobre o emprego formal são mantidas pelo Ministério do Trabalho e Emprego (MTE): a Relação Anual de Informações Sociais (RAIS) e o Cadastro Geral de Empregados e Desempregados (CAGED). Combinadas, elas cobrem os vínculos formais em todo o país, trazendo informações sobre admissões, desligamentos, salários, jornadas e ocupações classificadas pela Classificação Brasileira de Ocupações (CBO) \cite{RAIS2024, CAGED2024}.

\section{Contextualização e motivação}

Segundo \citeonline{Reis2022}, a engenharia de dados se firmou como uma disciplina central nas organizações modernas: é dela a responsabilidade de construir e manter os sistemas que transformam dados brutos em ativos úteis para a tomada de decisão. No âmbito público, a RAIS e o CAGED acumulam décadas de registros do mercado formal e formam um acervo de grande valor para análises de longo prazo.

O problema é que esses arquivos não foram pensados para consulta direta. Cada ano gera dezenas de gigabytes em formato texto, com \textit{layouts} e codificações que mudam ao longo da série, e os registros são só códigos, que precisam ser cruzados com tabelas de referência externas para fazer sentido \cite{RAIS2024, CAGED2024}. Sem uma camada de engenharia de dados que organize e padronize esse acervo, ele permanece subutilizado, e construir um \textit{pipeline} reproduzível que o transforme em séries temporais comparáveis é, por si só, uma contribuição relevante para quem trabalha com bases governamentais brasileiras.

Estudos sobre o mercado de TI brasileiro costumam trabalhar com recortes temporais curtos ou agregações setoriais publicadas por entidades como a BRASSCOM \cite{Brasscom2024} e a FECOMERCIOSP \cite{Fecomercio2024}. Olhando direto para os microdados da RAIS e do CAGED, é possível ir muito além: dá para quantificar o crescimento do setor, mapear disparidades regionais, medir o impacto de eventos macroeconômicos (como a recessão de 2015--2016, a pandemia de COVID-19 e a aceleração digital de 2021--2023) e projetar tendências a partir de quase duas décadas de registros formais.

\section{Problema de pesquisa}
\label{sec:problema}

A questão central que orienta o trabalho é: como construir uma plataforma analítica integrada que permita explorar e projetar o mercado formal de trabalho em TI no Brasil a partir dos microdados da RAIS e do CAGED, de forma reprodutível e com baixo custo operacional?

O objeto do trabalho é a plataforma: se dezenas de gigabytes de microdados governamentais heterogêneos podem ser tratados em uma infraestrutura modesta, de forma que outra pessoa consiga reproduzir o resultado. Uma plataforma analítica, porém, só se prova respondendo perguntas. Por isso ela foi posta à prova por uma pergunta sobre o mercado, \textit{o crescimento do emprego formal em tecnologia entre 2007 e 2026 veio acompanhado de melhores salários e de mais equidade?}, desdobrada nas cinco perguntas específicas do \autoref{quadro:perguntas}, que estruturam tanto o capítulo de resultados quanto as abas do \textit{dashboard}.

\begin{quadro}[htbp]
    \centering
    \caption{Perguntas específicas de pesquisa e a evidência usada para responder cada uma}
    \label{quadro:perguntas}
    \footnotesize
    \renewcommand{\arraystretch}{1.3}
    \begin{tabularx}{\textwidth}{>{\raggedright\arraybackslash}p{0.5cm}
                                 >{\raggedright\arraybackslash\hsize=1.05\hsize}X
                                 >{\raggedright\arraybackslash\hsize=0.95\hsize}X}
        \toprule
        \rowcolor{black!8}
        \textbf{\#} & \textbf{Pergunta} & \textbf{Evidência} \\
        \midrule
        1 & Quanto o mercado cresceu, e em que ritmo? & Estoque anual (RAIS) e saldo mensal (CAGED) \\
        2 & Onde está o trabalho de TI: em que empresas e em que lugares? & Cruzamento das lentes de setor e ocupação; agregações por unidade federativa \\
        3 & O emprego de TI melhorou em qualidade? & Remuneração em salários mínimos; giro anual; tábua de vida \\
        4 & Quem ganha menos, e por quê? & Decomposição de Oaxaca-Blinder por sexo e por raça/cor \\
        5 & Para onde o mercado vai? & Modelo SARIMA sobre o saldo mensal e estimador de \textit{nowcast} do estoque \\
        \bottomrule
    \end{tabularx}
    \\[4pt]
    Fonte: Elaborado pelo autor.
\end{quadro}

Nem as perguntas nem as fontes ficaram fixas desde a proposta: elas foram revistas a cada volta do ciclo do CRISP-DM (\autoref{sec:metodo}).

\section{Objetivos}
\label{sec:objetivos}

O objetivo geral foi desenvolver uma plataforma integrada para análise histórica, visualização interativa e previsão do mercado de trabalho em TI no Brasil, a partir dos microdados da RAIS e do CAGED, reprodutível e sem custo de infraestrutura.

Para alcançá-lo, foram definidos os seguintes objetivos específicos:

\begin{itemize}
    \item Implementar um \textit{pipeline} automatizado de engenharia de dados, em camadas, que processe os microdados da RAIS (2007--2025) e do CAGED (2007--2026) e se atualize quando o MTE publica dado novo;
    \item Definir e documentar um recorte de tecnologia que preserve as duas formas de delimitar o setor: pela atividade da empresa (CNAE) e pela ocupação do trabalhador (CBO);
    \item Traduzir os códigos pelos dicionários oficiais do MTE e auditar o resultado, verificando que a camada publicada pode ser reconstruída a partir do dado bruto;
    \item Desenvolver um \textit{dashboard} \textit{web} interativo, sem servidor de aplicação, organizado pelas perguntas de pesquisa;
    \item Prever o saldo mensal de empregos e estimar o estoque do ano corrente antes da publicação da RAIS, validando os dois contra estimadores ingênuos;
    \item Analisar a qualidade do emprego e as diferenças salariais por sexo e raça/cor;
    \item Publicar as camadas tratadas como dados abertos, para que a análise possa ser refeita por terceiros.
\end{itemize}

\section{Justificativa}

Do lado científico, a pesquisa oferece um retrato quantitativo do mercado de TI nacional cobrindo quase duas décadas, com detalhamento por ocupação (CBO), unidade federativa, setor econômico (CNAE) e porte de estabelecimento, um nível de granularidade que costuma ficar fora das análises publicadas sobre o setor. Esse tipo de evidência pode alimentar decisões na formulação de políticas de formação profissional, no planejamento de recursos humanos e em estudos de economia do trabalho.

No plano técnico, o trabalho testa se ferramentas leves e de código aberto, como DuckDB e Apache Parquet, substituem infraestruturas pesadas e caras em volumes dessa ordem \cite{Raasveldt2019, ApacheParquet2024}, e se um \textit{pipeline} desse porte pode rodar sem servidor próprio e sem custo de infraestrutura. Se puder, o mesmo desenho serve a outras bases governamentais brasileiras de grande volume.

\section{Estrutura do trabalho}

Este trabalho está organizado em cinco capítulos. O \autoref{cap:referencial} apresenta a fundamentação teórica: o mercado de TI e a literatura que orienta a análise, as bases de dados, a arquitetura de dados, os modelos de série temporal, os métodos de decomposição salarial e de sobrevivência e, por fim, os trabalhos relacionados. O \autoref{cap:metodologia} descreve os materiais e o método: as fontes, o recorte de tecnologia, a arquitetura implementada, as ferramentas e as revisões feitas em relação à proposta do TCC~1. O \autoref{cap:resultados} apresenta a plataforma, sua reprodutibilidade e os achados de qualidade do dado, e responde às cinco perguntas do \autoref{quadro:perguntas}. O \autoref{cap:consideracoes} responde à pergunta de pesquisa, confere cada objetivo, retoma as contribuições, declara as limitações e indica caminhos de continuidade.
